\documentclass{beamer}
% \usepackage[utf8]{inputenc}
\usepackage{default}
\mode<presentation>{
    %\usepackage{beamerthemeshadow}
    %\usepackage{beamerthemeshadow}
    %\usepackage{beamerthemePaloAlto}
    %\usepackage{beamerthemeSzeged}
    %\usepackage{beamerthemeAntibes}
    %\usepackage{beamerthemeBergen}
    %\usepackage{beamerthemeBerkeley}
    %\usepackage{beamerthemeBerlin}
    %\usepackage{beamerthemeBoadilla}  
    %\usepackage{beamerthemeboxes}  
    %\usepackage{beamerthemeCambridgeUS}
    %\usepackage{beamerthemeCopenhagen} 
    %\usepackage{beamerthemeDarmstadt}
    %\usepackage{beamerthemeDresden}
    %\usepackage{beamerthemeFrankfurt}
    %\usepackage{beamerthemeGoettingen}
    %\usepackage{beamerthemeHannover}
    \usepackage{beamerthemeIlmenau}
    %\usepackage{beamerthemeJuanLesPins} 
    %\usepackage{beamerthemeLuebeck} 
    \usepackage{pgf,pgfarrows,pgfnodes,pgfautomata,pgfheaps,pgfshade}
    \usepackage{clrscode3e}
    \beamertemplatetransparentcovereddynamic
    \beamertemplateballitem
    %\beamertemplatefootpagenumber
}
\mode<handout>{
    \usepackage[bar]{beamerthemetree}
    % Colocando um fundo cinza quando for gerar transparências para serem impressas
    % mais de uma transparência por página
    \beamertemplatesolidbackgroundcolor{black!5}
}
% \usepackage[ruled,chapter]{algorithm}
\usepackage{algorithm}	
\usepackage{algorithmic}

\usepackage{comment}
\usepackage{amsmath,amssymb}
\usepackage[brazil]{varioref}
\usepackage[english,brazil]{babel}
\usepackage[utf8]{inputenc}
\usepackage{graphicx}
\usepackage[alf]{abntex2cite}
\usepackage{abntex2abrev}
% \usepackage[round,sort,authoryear,nonamebreak]{natbib}
% \usepackage[authoryear]{natbib}
%\usepackage{minitoc-hyperref}
%\usepackage{listings}
%\usepackage{listings}
%\usepackage{colortbl}
%\usepackage{pslatex}
\beamertemplatetransparentcovereddynamic

\usepackage{tikz}
\usetikzlibrary{arrows}
\tikzstyle{block}=[draw opacity=0.7,line width=1.4cm]

\newcommand*\oldmacro{}%
\let\oldmacro\insertshorttitle%
\renewcommand*\insertshorttitle{%
  \oldmacro\hfill%
  \insertframenumber\,/\,\inserttotalframenumber}

\title[Trabalho de Conclusão de Curso]{Trabalho de Conclusão de Curso}
\subtitle{Desenvolvimento de um sistema web para a Associação dos Protetores da Bacia hidrográfica do Rio Gorutuba ``Kuruatuba''}

%\titlegraphic{}

\author[Guilherme Rocha Leite]{Guilherme Rocha Leite}
  \institute[UFVJM]{Universidade Federal dos Vales do Jequitinhonha e Mucuri \newline
  Bacharelado em Sistemas de Informação \newline
%   \inst{Departamento de Computação
	  
     Orientador: Prof.ª Erinaldo Barbosa da Silva\\
     Co-orientador: Thales Francisco Mota Carvalho\\
     $~$\\
%      Disciplina: Nome da Disciplina\\
}
\logo{\includegraphics[scale=0.1]{logo-ufvjm.jpg}}
% Se comentar a linha abaixo, irá aparecer a data quando foi compilada a apresentação
%\date{\textcolor{red}{III Encontro Mineiro de Equações Diferenciais, 2009}}
%\pgfdeclareimage[height=0.4cm]{das}{figs/logodas}
%\pgfdeclareimage[height=1cm]{logo}{logos}
% pode-se colocar o LOGO assim
%\logo{\pgfuseimage{logo}}
% ou...
%\logo{\vbox{\hbox to 3cm{\hfil\pgfuseimage{logo}}}}

\begin{document}
%\xdefinecolor{MyGreen}{rgb}{0,0.6,0}
%\beamerboxesdeclarecolorscheme{MyGreen}{MyGreen}{MyGreen!20!averagebackgroundcolor}
\beamerboxesdeclarecolorscheme{formula}{white}{blue!250!averagebackgroundcolor}
\frame{\titlepage}

%\frame{\titlepage}
\part{Presentation}
% \section{Sumário}

\frame{
\frametitle{Sumário}
\tableofcontents
}

\AtBeginSection[]{
  \frame<handout:0>{
    %\frametitle{Sumário}
    \tableofcontents[current,currentsection]
  }
}

%===================================Slide=================================================
		

\section{Motivação}
\begin{comment}
\begin{frame}
  \frametitle{Os custos logísticos nas organizações}
  \begin{itemize}
   \item Logística: pode ser definida como o processo de planejamento, implementação e controle dos fluxos de insumos e produtos na cadeia produtiva, de modo 
   que as mercadorias possam ser transportadas, desde as origens até os destinos, em tempo hábil e em conformidade com as necessidades de quem as demanda (CNT, 2014).
      
  \end{itemize}

\end{frame}
\end{comment}

\begin{frame}
  \frametitle{Os custos logísticos nas organizações}
  \begin{itemize}
   \item Os custos logísticos (estoque, transporte e armazenagem) representam grande parcela dos custos totais das empresas. \newline
   \item Sua maior parte se concentra no transporte, especialmente no Brasil: 59,8\% dos custos logísticos totais (CNT, 2014). \newline
   \item Fatores agravantes: custos do modal rodoviário, questões de legislação trabalhista e sanitária, bem como reutilização de materiais. \newline
  \end{itemize}
   
% Plano CNT de Transporte e Logística 2014  
% Brasil: 67% do total de cargas pelo rodoviário. EUA 31%
% problema grandes proporções em que pequenos percentuais de economia podem representar significativos ganhos, principalmente no Brasil
% A maior parte desses custos está relacionada ao transporte, especialmente no Brasil 59,8\% (CNT, 2014)
% Brasil, onde o modal rodoviário é o mais utilizado para transporte de cargas 
% É o modal mais caro no Brasil depois do aéreo
 
\end{frame}

\section{O Problema de Roteamento de Veículos}
\begin{frame}
 \frametitle{O Problema de Roteamento de Veículos}
 \begin{itemize}
  \item Deste modo, inúmeros esforços são empregados na tentativa de economizar qualquer precentual de custos de transporte    
  \item Entre as diversas pesquisas teóricas e práticas nesse campo, há O Problema de Roteamento de Veículos (VRP, do inglês \textit{Vehicle Routing Problem}) 
  \item O VRP consiste em definir rotas entre um depósito e um conjunto de pontos de entrega que minimize o custo de transporte (Pereira, 2007)
  \item Cada um destes pontos requer uma quantidade específica de produtos a serem entregues. Cada ponto de entrega é visitado apenas uma vez e por apenas um veículo, 
  o qual possui capacidade limitada (Caric e Gold, 2008b)
 \end{itemize}
  
\end{frame}


\begin{frame}
 \frametitle{O Problema de Roteamento de Veículos Multiobjetivo com Coleta Opcional - MOVRPOP}
 \begin{itemize}
  \item Há inúmeras variações do VRP original, as quais refletem as diversas situações encontradas na realidade.
  \item A variante aqui abordada é derivada do Problema de Roteamento de Veículos com Coleta e Entrega Simultâneas (VRPSPD) (Parragh et al., 2008).  
  \item No VRPSPD os consumidores apresentam demandas de \textbf{entrega e também de coleta}. Além disso, ambas as demandas devem ser \textbf{atendidas em uma única visita}. 
  \item Neste trabalho a obrigatoriedade do VRPSPD de atender as demandas de coleta é relaxada, \textbf{deixando de ser uma restrição e se tornando um objetivo a ser 
  minimizado}, assim como a distância total percorrida.
 \end{itemize}
  
\end{frame}

\begin{frame}
 \frametitle{MOVRPOP - Modelagem Matemática}
 \begin{itemize}
 \item $G = (V,A)$ é um grafo completo onde $V = \{v_0, v_1, v_2, \ldots, v_n\}$ é o conjunto de vétices, sendo $v_0$ o depósito e os demais vértices são consumidores;
 \item $A = \{(v_i,v_j) : i,j \in V ~ e ~ i \neq j\}$ é o conjunto de arestas, onde $(v_i,v_j)$ representa o caminho do vértice $v_i$ ao vértice $v_j$; 
 \item $c_{ij} \geq 0$ é um valor associado a cada aresta $A(v_i, v_j)$ e representa o custo para se alcançar o vértice $v_j$ a partir do vértice $v_i$;
 \item cada vértice $v_i$ possui uma demanda não-negativa de entrega ($d_i$) e coleta ($p_i$) associada a ele.
\end{itemize}
  
\end{frame}

\begin{frame}
 \frametitle{MOVRPOP - Modelagem Matemática}
 \begin{itemize}
 \item $k_{max}$ é o tamanho da frota de veículos disponíveis com capacidade $Q$;
 \item $y_{ij}$ é a variável de fluxo que representa o somatório das cargas coletadas entre o depósito e o nó $v_i$ (inclusive), dirigido ao nó $v_j$, 
 \item $z_{ij}$ é a variável de fluxo que representa o somatório das cargas entregues entre o nó $v_i$ (exclusive) e o depósito, dirigido ao nó $v_j$
\end{itemize}
  
\end{frame}

\begin{frame}
 \frametitle{MOVRPOP - Modelagem Matemática}
 Variáveis de Decisão do problema:
 
 {\small
\begin{equation*}
 x^{k}_{ij} = \begin{cases}
               1 \qquad \textrm{, se a aresta $(v_i,v_j)$ faz parte da rota trafegada pelo veículo $k$,}\\
               0 \qquad \textrm{, caso contrário.}
              \end{cases}
\end{equation*}
\begin{equation*}
 \ell_{j} = \begin{cases}
           1 \qquad \textrm{, se a demanda de coleta do consumidor $v_j$ é satisfeita,}\\
           0 \qquad \textrm{, caso contrário.}
          \end{cases}
\end{equation*}
}
  
\end{frame}

\begin{frame}
 \frametitle{MOVRPOP - Modelagem Matemática}
 Funções-objetivo: \newline \newline
 - Minimização do custo
 \begin{equation}
\label{eq:modelagem_1}
Min ~f(\textbf{x}) ~=~ \sum^{k_{max}}_{k=1}\sum^{n}_{i=0}\sum^{n}_{j=0}c_{ij}x^{k}_{ij}
\end{equation}
%

 - Minimização das coletas não realizadas
\begin{equation}
\label{eq:modelagem_2}
Min ~f(\boldsymbol{\ell}) ~=~  \sum^{n}_{j=1}p_{j}(1-\ell_{j})
\end{equation}
  
\end{frame}

\begin{frame}
 \frametitle{MOVRPOP - Modelagem Matemática}
 \begin{itemize}
 \item cada vértice deve ser visitado por um único veículo:
% cada ponto deve ser visitado por um único veículo
\begin{equation}
\label{eq:modelagem_3}
\sum^{n}_{i=0}\sum^{k_{max}}_{k=1}x^{k}_{ij} = 1, \qquad j = 1,\ldots,n
\end{equation}

\item a conservação do fluxo deve ser garantida: 
% conservação de fluxo
\begin{equation}
\label{eq:modelagem_4}
\sum^{n}_{i=0}x^{k}_{ij} - \sum^{n}_{i=0}x^{k}_{ji} = 0, \qquad j = 0,\ldots,n \textrm{ e } k = 0,\ldots,k_{max}
\end{equation}

\item no máximo $k_{max}$ veículos podems ser utilizados:
% k veículos, no máximo devem ser utilizados
\begin{equation}
\label{eq:modelagem_5}
\sum^{n}_{j=1}x^{k}_{0j} \leq 1, \qquad k = 1,\ldots,k_{max}
\end{equation}
  \end{itemize}

\end{frame}

\begin{frame}
 \frametitle{MOVRPOP - Modelagem Matemática}
 \begin{itemize}
 \item todas as demandas de entrega devem ser satisfeitas:
% demandas de entrega devem ser satisfeitas em cada consumidor
\begin{equation}
\label{eq:modelagem_6}
\sum^{n}_{i=0}z_{ij} - \sum^{n}_{i=0}z_{ji} = d_{j}, \qquad \forall j \neq 0
\end{equation}

\item uma vez que a variável $l_j$ assumir valor 1, a demanda de coleta do consumidor $v_j$ deve ser satisfeita:
% demandas de coleta devem ser satisfeitas em cada consumidor
\begin{equation}
\label{eq:modelagem_61}
\sum^{n}_{i=0}y_{ji} - \sum^{n}_{i=0}y_{ij} = p_j \ell_j, \qquad \forall j \neq 0
\end{equation}

\item os veículos possuem uma capacidade limite para o transporte de cargas. Esse transporte deve ocorrer através das arestas incluídas na solução:
% demandas devem ser transportadas nos arcos incluidos na solução
% e ainda impõe um limite para a carga total transportada por veículo
\begin{equation}
\label{eq:modelagem_7}
y_{ij} + z_{ij} \leq Q \sum^{k_{max}}_{k=1}x^{k}_{ij}, \qquad i, j = 0,\ldots,n
\end{equation}
  \end{itemize}

\end{frame}

\begin{frame}
 \frametitle{MOVRPOP - Modelagem Matemática}
 \begin{itemize}
 \item restrição de integralidade das variáveis de decisão:
% restrição de integralidade
\begin{equation}
\label{eq:modelagem_8}
x_{ij}, \ell_j \in \{0, 1\}, \qquad i, j = 0,\ldots,n
\end{equation}

\item restrição de não-negatividade para fluxo de demandas de coleta:
% restrição de não negatividade
\begin{equation}
\label{eq:modelagem_9}
y_{ij} \geq 0, \qquad i, j = 0,\ldots,n
\end{equation}

\item restrição de não-negatividade para fluxo de demandas de entrega:
% restrição de não negatividade
\begin{equation}
\label{eq:modelagem_10}
z_{ij} \geq 0, \qquad i, j = 0,\ldots,n
\end{equation}
  \end{itemize}

\end{frame}

\section{Busca Local Multiobjetivo Baseada em Indicador}
\begin{frame}
 \frametitle{Características do IBMOLS}
 \begin{itemize}
  \item Nem sempre há uma \textbf{clara definição do que seja uma boa aproximação} da Fronteira de Pareto    
  \item A qualidade da solução dependerá do contexto do problema e das \textbf{preferências do tomador de decisão} 
  \item Embora os algoritmos evolucionários sejam os mais utilizados, é reconhecida a eficiência dos \textbf{métodos 
  baseados em busca local}, principalmente em problemas de \textbf{otimização combinatória de larga escala} (Basseur, 2007)
  
 \end{itemize}
  
\end{frame}

\begin{frame}
 \frametitle{Características do IBMOLS}
 \begin{itemize}
  \item Isto posto, o presente trabalho utiliza para resolução do MOVRPOP a busca local multiobjetivo baseada em indicador (IBMOLS, do inglês \emph{Indicator-based 
  multiobjetive local search}).    
  \item No IBMOLS é utilizado um indicador binário (${I}$) que expressa a preferência do tomador de decisão e \textbf{atribui a cada solução um valor} real que reflete 
  a sua qualidade
  \item A seguir é apresentado a versão adotada do algoritmo do IBMOLS, proposta por (Assis, 2013)
  
 \end{itemize}
  
\end{frame}

\begin{frame}
%  \frametitle{Algoritmo IBMOLS}
%\begin{algorithm}
\tiny{
 \begin{algorithmic}[1]
  \REQUIRE N (Tamanho da População), I(Indicador Binário) %1
  \STATE $PO \leftarrow \emptyset$; %2
  \REPEAT %3
  		\STATE $P \leftarrow $ População Inicial($N$);  %4
  		\STATE $Aptidao[~] \leftarrow $Calcular Aptidão($P$,$I$) %5
  		\STATE $s_{w} \leftarrow$ Menor($Aptidao$); %6
  		\FORALL{solução $s \in P$} %7
  				\STATE continue $\leftarrow$ TRUE; %8
				\WHILE {continue}  				%9
  				\STATE $s' \leftarrow$ Perturbação($s$);  %10
  				\STATE $s'' \leftarrow$ Busca Local($s'$); %11
   				\IF { $s= s''$} %12
   					\STATE $s'' \leftarrow s'$;  continue $\leftarrow$ FALSE; %13
   				\ENDIF %14
   				\IF{$Aptidao(s'') > Aptidao(s_{w})$} 
   					\STATE $P \leftarrow P \cup \{s''\} - \{s_{w}\}$;
   					\STATE Calcular Aptidão($s''$, $s_{w}$, $P$, $I$);
   					\STATE $s_{w} \leftarrow$ Menor($Aptidao$); 
   				\ELSE 
   					\STATE continue $\leftarrow$ FALSE;
   				\ENDIF
   				\ENDWHILE
  		\ENDFOR
  		 \STATE $PO \leftarrow$ Soluções Não Dominadas$(PO \cup P)$;
  \UNTIL {Critério de Parada}
  \STATE \textbf{retornar} $PO$;
 \end{algorithmic} }
%   \caption{Algoritmo IBMOLS modificado}
% \label{alg:ibmolsassis}
%\end{algorithm}  

\end{frame}

\begin{frame}
 \frametitle{Indicadores Binários}
 \begin{itemize}
  \item A qualidade das soluções no $I_{\epsilon}$ é mensurada pela a distância entre duas soluções no espaço dos objetivos. Seu cálculo é efetuado da seguinte maneira: 
 \begin{equation}
 \label{eq:indicadorepsilon}
 I_{\epsilon}(x_1, x_2) = max_{i \in \{1,...,n\}}(f_i(x_1)-f_i(x_2))
 \end{equation}
  \item O cálculo de aptidão de um indivíduo é determinado pelo somatório (Basseur, 2007):
  \begin{equation}
\label{eq:indicador3}
I(P\setminus\{x\},x) = \sum_{z \in P \setminus \{x\}} -e^{-I_{\epsilon}(z,x)/k}
\end{equation}
 \end{itemize}
  
\end{frame}

\begin{frame}
 \frametitle{Indicadores Binários}
 \begin{itemize}
  \item Dadas duas soluções $x_1$ e $x_2$, o indicador binário de hipervolume $I_{HD}(x_1,x_2)$, representa o volume do espaço de objetivos que é dominado por $x_1$ mas 
  não é dominado por $x_2$: 
 \begin{equation}
 \label{eq:indicadorhd}
 I_{HD}(x_1, x_2) = \begin{cases}
          $H${(x_1)}$ - H$(x_2)$$ \qquad \textrm{ se $x_1$ $\prec$ $x_2$},\\
          $H$(x_1+x_2)$ - H$(x_2)$$ \qquad \textrm{ caso contrário.}
             \end{cases}
\end{equation}
  \item O cálculo de aptidão de um indivíduo é determinado pelo somatório:
  \begin{equation}
\label{eq:indicador3}
I(P\setminus\{x\},x) = \sum_{z \in P \setminus \{x\}} -e^{-I_{HD}(z,x)/k}
\end{equation}
 \end{itemize}
  
\end{frame}

\begin{frame}
 \frametitle{Indicadores Binários}
 \begin{itemize}
  \item O indicador $I_{FON}$ basicamente propõe um ranqueamento entre as soluções, a partir do número de soluções que cada uma domina: 
\begin{equation}
 \label{eq:indicadorfon}
 I_{Fon}(x_1,x_2) =  \begin{cases}
          -1 \qquad \textrm{ se $x_1$ $\prec$ $x_2$},\\
          0 \qquad \textrm{ caso contrário.}
             \end{cases}
\end{equation}
  \item Direfentemente dos outros indicadores, o cálculo de aptidão dos indivíduos da população é realizado através do simples somatório dos valores de fitness de cada 
  indivíduo (Fonseca et al. (1993)):
  \begin{equation}
\label{eq:indicador1}
I(P\setminus\{x\},x) = \sum_{z \in P \setminus \{x\}} I_{FON}(z,x)
\end{equation}

 \end{itemize}
  
\end{frame}

% Amanhã: métricas utilizadas, parâmetros empregados, tabelas coloridas, conclusão
\section{Testes Computacionais}

\begin{frame}
  \frametitle{Instâncias de Problemas e Parâmetros utilizados}
    
    \begin{itemize}
    \item Foram utilizadas 12 instâncias de 50 a 150 consumidores mais o depósito, propostas por Salhi e Nagy (1999) para o VRPSPD.
    \item Cada instância foi executada 5 vezes com cada indicador
    \item O número de gerações do IBMOLS foi 50, tamanho da população inicial 10, taxa de mutação dos mecanismos de perturbação 0.3 e o fator de escala $k$ (no caso dos 
    indicadores $I_{\epsilon}$ e $I_{HD}$) 0.1. 
    \end{itemize}
  
\end{frame}

\begin{frame}
  \frametitle{Métricas utilizadas para avaliação}
    Foram utilizadas as seguintes métricas na avaliação do desempenho dos algoritmos:
    \begin{itemize}
    \item Tempo de execução
    \item Cardinalidade, que mensura quantas possibilidades de escolha o tomador de decisão pode ter
    \item Hipervolume: a qualidade de uma fronteira é definida pela soma dos hipercubos de cada uma das soluções que a compõem (Hughes, 2005)
    \item Cobertura: determina em que medida uma fronteira de Pareto domina as demais
    \end{itemize}
  
\end{frame}

\begin{frame}
  \frametitle{Resultados de Tempo de Execução}
    \footnotesize{
    \begin{tabular}{c|ccc}
\hline
Instância & \textcolor{red}{$I_{\epsilon}$}             & \textcolor{green}{$I_{FON}$}             & \textcolor{blue}{ $I_{HD}$}      \\
\hline
\textcolor{red}{CMT1X}     & 4783,25 (241,31)    & 5221,07 (448,8)     & 5140,48 (202,24)    \\
\textcolor{green}{CMT1Y}     & 7700,22 (234,36)    & 7246,09 (964,65)    & 7366,17 (477,8)     \\
\textcolor{red}{CMT2X}     & 11696,6 (1009,23)   & 11868,65 (1620,46)  & 11772,4 (649,95)    \\
\textcolor{red}{CMT2Y}     & 27499,58 (2073,65)  & 32398,84 (1607,76)  & 32438,68 (4541,82)  \\
\textcolor{red}{CMT3X}     & 152092,6 (6979,93)  & 162727,6 (11602,19) & 160291,6 (7381,23)  \\
\textcolor{red}{CMT3Y}     & 203925,6 (7675,61)  & 207824,2 (10057,66) & 207873,4 (4173,17)  \\
\textcolor{red}{CMT4X}     & 574257,6 (18919,79) & 588513,2 (16569,53) & 574362,2 (37086,5)  \\
\textcolor{red}{CMT4Y}     & 784814,4 (45346,66) & 801915,6 (16155,43) & 790977 (57675,27)   \\	
\textcolor{red}{CMT11X}    & 497025,4 (38917,17) & 506154,8 (44709,15) & 507026 (52041,12)   \\
\textcolor{red}{CMT11Y}    & 327682,2 (34860,79) & 345287,6 (19075,34) & 342428,4 (18764,02) \\
\textcolor{red}{CMT12X}    & 100029,28 (4513,96) & 105918,5 (5543,89)  & 105928,4 (4785,45)  \\
\textcolor{blue}{CMT12Y}    & 129340,2 (6356,13)  & 138568,2 (5581,87)  & 128980,2 (5395,27)   \\
\hline
\textcolor{red}{Média}& 235070,58 (13927,38)& 242803,70 (11161,39)& 239548,74 (16097,82)\\
\hline
\end{tabular}
  }
\end{frame}


\begin{frame}
  \frametitle{Resultados de Cardinalidade}
    \footnotesize{
    \begin{tabular}{c|ccc}
\hline
Instância & \textcolor{red}{$I_{\epsilon}$}             & \textcolor{green}{$I_{FON}$}             & \textcolor{blue}{ $I_{HD}$}      \\
\hline
\textcolor{blue}{ CMT1X}     & 2,6 (0,49)  & 2,8 (0,4)   & 3,6 (0,8)   \\
\textcolor{green}{CMT1Y}     & 7,4 (1,02)  & 11,4 (1,02) & 11,2 (1,17) \\
\textcolor{red}{CMT2X}     & 5 (1,41)    & 4,4 (1,62)  & 3,2 (0,75)  \\
\textcolor{green}{CMT2Y}     & 22,4 (3,01) & 29,8 (3,06) & 27,8 (5,08) \\
\textcolor{blue}{ CMT3X}     & 5,2 (1,47)  & 6,2 (2,14)  & 7 (1,1)     \\
\textcolor{green}{CMT3Y}     & 17,2 (1,6)  & 19,8 (3,12) & 18 (3,03)   \\
\textcolor{green}{CMT4X}     & 7 (1,26)    & 7,2 (1,94)  & 6,4 (1,5)   \\
\textcolor{red}{CMT4Y}     & 30,8 (3,06) & 29 (3,1)    & 30,2 (2,93) \\
\textcolor{red}{CMT11X}    & 37,4 (7,66) & 35,6 (6,89) & 37 (8,37)   \\
\textcolor{blue}{CMT11Y}    & 20,8 (1,6)  & 24,6 (4,08) & 24,8 (2,93) \\
\textcolor{green}{CMT12X}    & 14,6 (2,73) & 16 (3,58)   & 14,8 (2,99) \\
\textcolor{red}{CMT12Y}    & 27,2 (4,45) & 24,8 (4,02) & 24,2 (3,87) \\
\hline
\textcolor{green}{Média}& 16,47 (2,48)& 17,63 (2,91)& 17,35 (2,88)\\
\hline
\end{tabular}
  }
\end{frame}

\begin{frame}
  \frametitle{Resultados de Hipervolume}
    \footnotesize{
    \begin{tabular}{c|ccc}
\hline
Instância & \textcolor{red}{$I_{\epsilon}$}             & \textcolor{green}{$I_{FON}$}             & \textcolor{blue}{ $I_{HD}$}      \\
\hline
\textcolor{blue}{CMT1X}     & 391,61 (16,62)     & 405,05 (17,78)     & 413,48 (3,61)      \\
\textcolor{green}{CMT1Y}     & 2833,78 (244)      & 3119,52 (80,61)    & 3016,16 (107,51)   \\
\textcolor{red}{CMT2X}     & 2166,34 (247,25)   & 1886,7 (351,62)    & 1956,23 (266,36)   \\
\textcolor{red}{CMT2Y}     & 54779,43 (745,53)  & 54355,08 (837,48)  & 54242,71 (1371,66) \\
\textcolor{green}{CMT3X}     & 970,92 (155,08)    & 1123,73 (26,12)    & 1116,95 (5,63)     \\
\textcolor{green}{CMT3Y}     & 9247,71 (93,23)    & 9348,28 (40,43)    & 9106,29 (219,12)   \\
\textcolor{green}{CMT4X}     & 1890,53 (285,22)   & 2296,68 (634,4)    & 1820,75 (348,09)   \\
\textcolor{red}{CMT4Y}     & 38618,37 (866,12)  & 38369,73 (1103,39) & 38393,53 (1582,99) \\
\textcolor{red}{CMT11X}    & 36891,88 (1131,33) & 36843,05 (614,28)  & 36393,13 (912,33)  \\
\textcolor{red}{CMT11Y}    & 19246,31 (1584,18) & 18610,46 (1585,78) & 18651,36 (1355,23) \\
\textcolor{green}{CMT12X }   & 10735 (70,94)      & 10736,53 (98,45)   & 10616,73 (222,36)  \\
\textcolor{green}{CMT12Y}    & 20577,81 (1193,77) & 20770,49 (535,56)  & 20574,07 (871,64)  \\
\hline
\textcolor{red}{Média} & 16529,14 (552,77)  & 16488,78 (493,83)  & 16358,45 (605,54)  \\
\hline
\end{tabular}
  }
\end{frame}

\begin{frame}
  \frametitle{Resultados de Cobertura}
    \footnotesize{
    \begin{tabular}{c|ccc}
\hline
Instância & \textcolor{red}{$I_{\epsilon}$}             & \textcolor{green}{$I_{FON}$}             & \textcolor{blue}{ $I_{HD}$}      \\
\hline
\textcolor{green}{CMT 1X}  & 0,68 (0,19) & 0,74 (0,25) & 0,73 (0,22) \\
\textcolor{blue}{CMT 1Y}  & 0,32 (0,09) & 0,50 (0,08) & 0,51 (0,18) \\
\textcolor{blue}{CMT 2X}  & 0,46 (0,45) & 0,36 (0,34) & 0,46 (0,20) \\
CMT 2Y  & 0,50 (0,16) & 0,34 (0,15) & 0,50 (0,09) \\
\textcolor{blue}{CMT 3X}  & 0,30 (0,19) & 0,48 (0,12) & 0,62 (0,12) \\
\textcolor{green}{CMT 3Y}  & 0,36 (0,08) & 0,52 (0,17) & 0,37 (0,13) \\
\textcolor{green}{CMT 4X}  & 0,32 (0,14) & 0,64 (0,26) & 0,37 (0,23) \\
\textcolor{blue}{CMT 4Y}  & 0,42 (0,16) & 0,36 (0,10) & 0,48 (0,11) \\
\textcolor{red}{CMT 11X} & 0,50 (0,17) & 0,45 (0,12) & 0,44 (0,27) \\
\textcolor{blue}{CMT 11Y} & 0,40 (0,21) & 0,37 (0,12) & 0,40 (0,13) \\
\textcolor{blue}{CMT 12X} & 0,45 (0,15) & 0,40 (0,03) & 0,51 (0,17) \\
\textcolor{blue}{CMT 12Y}& 0,37 (0,11) & 0,31 (0,08) & 0,48 (0,10) \\
\hline
\textcolor{blue}{Média} & 0,42 (0,18) & 0,46 (0,15) & 0,49 (0,16)\\
\hline
\end{tabular}
  }
\end{frame}

\section{Conclusão}

\begin{frame}
  \frametitle{Conclusão}
    
    \begin{itemize}
    \item Os custos de transporte representam lugar de destaque nas preocupações das organizações. Qualquer percentual de economia pode representar um grande avanço
    \item A revisão bibliográfica mostrou que há poucos trabalhos de VRP que tratam de demandas de coleta e entrega numa abordagem multiobjetivo.
    \item Por ser um problema de complexidade $\mathcal{NP}$-difícil, busca-se o emprego de técnicas heurísticas para obtenção de soluções viáveis em um tempo 
    computacional aceitável
    \end{itemize}
  
\end{frame}

\begin{frame}
  \frametitle{Conclusão}
    
    \begin{itemize}
    \item A técnica empregada para a resolução do MOVRPOP foi a busca local local baseada em indicador (IBMOLS), que possui um indicador através do qual o tomador de 
    decisão pode explicitar suas preferências de otimização
    \item O IBMOLS foi testado com três diferentes tipos de indicadores encontrados na literatura $I_{\epsilon}$, $I_{HD}$ e o $I_{FON}$
    \item Os resultados obtidos demonstram um desempenho muito próximo entre os indicadores, sendo necessário um maior número de testes para se verificar se as diferenças
    entre eles são significativas do ponto de vista estatístico
    \end{itemize}
  
\end{frame}

\section{Trabalhos Futuros}
\begin{frame}
  \frametitle{Trabalhos Futuros}
    
    \begin{itemize}
    \item executar um maior número de testes para uma análise estatística melhor das diferenças entre os algoritmos;
    \item testar o IBMOLS em outros conjuntos de instâncias da literatura bem como em outros problemas de otimização;
    \item comparar o desempenho do IBMOLS em relação a outras meta-heurísticas;
    \item implementar e testar os indicadores utilizados em algoritmos evolucionários;
    \item testar o IBMOLS com outros indicadores encontrados na literatura e, a partir da comparação entre os mesmos através da observação e experiência adquiridas, 
    propor um novo indicador;
    \item aplicar a técnica proposta em um problema real.
    \end{itemize}
  
\end{frame}

\begin{frame}
    \cite{gonccalves2005proposta}
        
    \cite{de2003educaccao}

    \cite{Kuruatuba2011}

    \cite{elias}

    \cite{Sommerville}
\end{frame}

    

\section{Referências Bibliográficas}

\begin{frame}
    \frametitle{Referências Bibliográficas}
    \scriptsize{  
        \bibliographystyle{abnt-alf}
        \bibliography{referencias}
    }
  
\end{frame}

\end{document}
